\chapter[Limitations and Ethical Considerations]{Limitations and\linebreak Ethical Considerations}\label{ch:limits}

\section{Limitations}
The results should be read within the scope of the setup that produced them.
\begin{itemize}
\item \textbf{Scope.} Behavioural probes were exercised on English toxicity classification with Civil Comments, using local \code{transformers} sequence classification. Other modalities and tasks are not covered, and SHAP and LIME are not implemented.
\item \textbf{Probes and mapping.} Robustness is a random character swap, Explainability reads the model card, and each dimension yields one O/S/D triple. The O/S/D mapping is a proposal, scaled to the injected defects and not validated on fresh data.
\item \textbf{Experiments.} Sweeps use four or five levels, one model family and one training seed per condition. Defects are synthetic, and three of the nine suite models carry documentation-only defects. Fairness conclusions are specific to the binary ``identity mentioned'' attribute and the 3\,000-row sample.
\item \textbf{LLM interpretation.} \code{llm\_v1} depends on third-party providers with quotas, so part of the LLM-rated data fell back to the heuristic. It is kept as an experimental comparison engine.
\item \textbf{Integrity.} A hash match shows the bytes equal the manifest, not that the weights are free of backdoors.
\item \textbf{Reviewer study.} Seven submissions on twelve author-chosen packets, six of them under the shared code \code{R0}. The study shows where an independent reading differs from TrustLens and is not a statistical validation.
\item \textbf{Deployment.} Authentication was removed (migration \code{008\_remove\_auth\_and\_ownership}); the system is intended for local, single-user use.
\end{itemize}
Threats to validity follow the same pattern. Each probe measures a proxy rather than ``trustworthiness'' itself, post-hoc decisions are listed in Table~\ref{tab:posthoc}, and the small $n$ rules out significance tests.

\section{Ethical considerations}
\paragraph{Data.} The experiments use the public Civil Comments dataset \cite{borkan2019}, which contains offensive text and annotator identity labels. We used the labels only to define a coarse group flag for measuring group differences in a classifier, not to infer anything about individuals. Reviewers saw card text and numbers, not comment text.

\paragraph{Participants.} Reviewers used an assigned code, with no name or email, and no demographic data were collected. The raw export is kept unmodified and its hash recorded. The repository documents no formal institutional ethics review; this should be closed before a larger panel is recruited.

\paragraph{External services.} \code{llm\_v1} sends model-card text and probe evidence to third-party LLM providers, so auditing a private model with it would disclose that material. The default deterministic engine sends nothing. Account identifiers in provider errors were replaced with \code{org\_<redacted>}.

\paragraph{Licences and models.} Models are used under their own licences. Fine-tuned experimental models are research artifacts and are not published as moderation systems.

\paragraph{Responsible use.} A dashboard score can be read as a certificate, so the profile is the primary output, abstention is visible, and every proposed rating is labelled ``PROPOSED / REQUIRES VALIDATION''. TrustLens is meant to help a human reviewer prioritise questions: where the evidence is thin, where a measured value differs from a reference, where a card contradicts itself. It should not be used to approve, certify or rank models, or to make deployment decisions without review.

\paragraph{Resources.} Everything ran on one laptop GPU, and fine-tuning each model took under three minutes.
